\section{Rzuty kostką}

\subsection{Ogólna zasada}
Gracze często będą musieli rzucać jedną lub kilkoma kostkami. Ich liczba oraz rodzaj określony będzie przy pomocy kodu.

\subsection{Kod rzutu kostką}
Kod składa się z cyfry, po której następuje litera K, a po niej następna cyfra.
Czasem do tego zasadniczego członu kodu dodawana jest jeszcze jedna cyfra, np. 1K6+1.
Pierwsza cyfra określa liczbę kostek, którymi należy rzucić, litera K oznacza po prostu kostkę, zaś następująca po niej cyfra~--- rodzaj kostki lub kostek (trój- lub sześcienne).
Jeżeli wymagany jest rzut kostką trójścienną, należy rzucić normalną kostką o sześciu ścianach, a rezultat podzielić przed 2~--- ułamki zaokrąglając w górę.
Trzecia cyfra jest dodawana do lub odejmowana od liczby wyrzuconych oczek.
Jeżeli trzecia cyfra jest poprzedzona znakiem x, to liczbę wyrzuconych oczek należy przez tę cyfrę pomnożyć.

Przykład: 1K6x4 oznacza, że należy rzucić jedną kostką sześcienną, a rezultat pomnożyć przez 4.
1K3+2 oznacza, że należy rzucić jedną kostką sześcienną, rezultat podzielić przez 2 i do uzystkanego wyniku dodać 2.

\subsection{Skrócony kod rzutu kostką}
Na żetonach reprezentujących potwory powyższy kod podany jest w formie skróconej. Pierwsza cyfra jest liczbą rzucanych kostek sześciennych (zawsze!), zaś druga powinna być dodana do liczby wyrzuconych oczek, np. 1+1 oznacza to samo co 1K6+1.
